\documentclass{article}

\usepackage{amsmath,amssymb}
\usepackage{xurl}
\usepackage[hidelinks]{hyperref}
\setlength{\emergencystretch}{2em}

\input{command}

\title{The Abstract Schwarz-Map Hypotheses in Wolf Theorem 6.12}
\author{}
\date{2026-07-18}

\begin{document}
\maketitle
\gapnote{scope-restriction}{open}

\section{Source statement}

A Schwarz map in Chapter~5 of Ref.~\cite{Wolf2012Quantum} is a positive unital
matrix map $E$ satisfying source Equation~(5.2):
\begin{align}
  E(A^\dagger A)-E(A^\dagger)E(A)
    &\succeq0
    \qquad\text{for every }A.
  \label{eq:wolf_thm6_12_abstract_schwarz_fixed_points_schwarz_inequality}
\end{align}
This definition appears immediately before Example~5.3, at local source
lines~347--352.

Theorem~6.12 of Ref.~\cite{Wolf2012Quantum} assumes a unital Schwarz map and
an arbitrary full-rank fixed point of its trace-pairing adjoint. Its algebraic
conclusion is
\begin{align}
  \operatorname{Fix}(E)
    &=\{X:E(X)=X\},
  \label{eq:wolf_thm6_12_abstract_schwarz_fixed_points_fixed_space_definition}\\
  \operatorname{Fix}(E)
    &\text{ is a unital $*$-subalgebra}.
  \label{eq:wolf_thm6_12_abstract_schwarz_fixed_points_source_algebra_conclusion}
\end{align}
The proof is at local source lines~1395--1417. It invokes Proposition~6.9 on
positive fixed points, at local source lines~1161--1192, to replace the
full-rank witness by a positive-definite fixed point before applying the
weighted-trace argument. The complete source theorem also invokes the explicit
block realization in source Equation~(1.39) of Ref.~\cite{Wolf2012Quantum}. This note concerns only the
$*$-algebra conclusion in
(\ref{eq:wolf_thm6_12_abstract_schwarz_fixed_points_source_algebra_conclusion}).

The source theorem does not repeat the word ``positive'' because positivity is
part of the Chapter~5 definition of a Schwarz map. Making positivity explicit
in a formal signature does not strengthen the theorem.

\section{Weighted Schwarz equality}

Assume that $\rho>0$ satisfies
\begin{align}
  E^*(\rho)
    &=\rho,
  \label{eq:wolf_thm6_12_abstract_schwarz_fixed_points_invariant_weight}
\end{align}
and let $X\in\operatorname{Fix}(E)$. Positivity of $E$ implies
\begin{align}
  E(X^\dagger)
    &=X^\dagger.
  \label{eq:wolf_thm6_12_abstract_schwarz_fixed_points_adjoint_fixed}
\end{align}
Define the Schwarz gap by
\begin{align}
  \Delta_X
    &=E(X^\dagger X)-E(X^\dagger)E(X),
  \label{eq:wolf_thm6_12_abstract_schwarz_fixed_points_schwarz_gap_definition}\\
  \Delta_X
    &\succeq0.
  \label{eq:wolf_thm6_12_abstract_schwarz_fixed_points_schwarz_gap_positive}
\end{align}
The trace-pairing identity gives the calculation in source Equation~(6.60) of Ref.~\cite{Wolf2012Quantum}:
\begin{align}
  \tr(\rho\Delta_X)
    &=\tr\!\left(\rho E(X^\dagger X)\right)
      -\tr\!\left(\rho E(X^\dagger)E(X)\right),
  \label{eq:wolf_thm6_12_abstract_schwarz_fixed_points_weighted_gap_expand}\\
    &=\tr\!\left(E^*(\rho)X^\dagger X\right)
      -\tr\!\left(\rho X^\dagger X\right),
  \label{eq:wolf_thm6_12_abstract_schwarz_fixed_points_trace_adjoint_step}\\
    &=0.
  \label{eq:wolf_thm6_12_abstract_schwarz_fixed_points_weighted_gap_zero}
\end{align}
Because $\rho$ is positive definite and $\Delta_X$ is positive semidefinite,
(\ref{eq:wolf_thm6_12_abstract_schwarz_fixed_points_weighted_gap_zero})
implies
\begin{align}
  \Delta_X
    &=0.
  \label{eq:wolf_thm6_12_abstract_schwarz_fixed_points_schwarz_equality}
\end{align}
Repeating the argument for $X^\dagger$ gives equality in both one-sided
Schwarz inequalities. The abstract multiplicative-domain theorem then yields,
for fixed points $X$ and $Y$,
\begin{align}
  E(XY)
    &=E(X)E(Y),
  \label{eq:wolf_thm6_12_abstract_schwarz_fixed_points_multiplicativity}\\
    &=XY.
  \label{eq:wolf_thm6_12_abstract_schwarz_fixed_points_product_fixed}
\end{align}
This proves closure under multiplication; the preceding identities give
closure under adjoints, and the unital hypothesis supplies the identity.

The argument is formalized for arbitrary positive unital Schwarz maps once a
positive-definite invariant weight is supplied explicitly.\footnote{Formal
declarations:
\leanid{SchwarzMap.schwarz_equality_of_mem_fixedPoints},
\leanid{SchwarzMap.mem_multiplicativeDomain_of_mem_fixedPoints}, and
\leanid{SchwarzMap.fixedPointsStarSubalgebra} in
\doclink{QICLean/Channel/FixedPoint/AbstractAlgebra}.}

\section{Comparison with the Kraus specialization}

The earlier declaration
\leanid{Kraus.fixedPointsStarSubalgebra} assumed a Kraus representation and is
located in
\doclink{QICLean/Channel/FixedPoint/Algebra}. It is mathematically correct, but
complete positivity is stronger than the Schwarz-map hypothesis of
Theorem~6.12 in Ref.~\cite{Wolf2012Quantum}. The declaration remains a useful
specialization of the post-reduction lemma, but it is not the literal
formalization of the source theorem.

\section{Verdict}

The earlier Kraus-only attribution was a \emph{scope restriction}, and the
abstract lemma removes that restriction. A second scope restriction remains.
The formal signature assumes the positive-definite invariant weight in
(\ref{eq:wolf_thm6_12_abstract_schwarz_fixed_points_invariant_weight}), whereas
Theorem~6.12 assumes only a full-rank fixed point and derives a faithful
positive witness through Proposition~6.9
\cite[Theorem~6.12 and Proposition~6.9]{Wolf2012Quantum}. The current theorem
is therefore the source-faithful post-reduction algebra lemma, not yet a
literal formalization of Theorem~6.12.

Eliminating the remaining restriction requires formalizing the proposition on
positive fixed parts for the abstract positive trace-preserving trace adjoint,
deriving a positive-definite fixed point from the full-rank witness, and then
applying the present lemma. The explicit trace-adjoint realization from source
Equation~(1.39) of Ref.~\cite{Wolf2012Quantum} also remains open at this level of generality. The downstream
Schr{\"o}dinger-picture density-block classification, support reduction, and
zero-block assembly required for Theorem~6.14 are complete; see
\path{docs/paper-gaps/wolf_theorem6_14_fixed_point_projection_gap.tex}.

\bibliographystyle{alpha}
\bibliography{references}

\end{document}
